% Created 2026-11-04 Wed 21:53
% Intended LaTeX compiler: pdflatex
\documentclass[11pt]{article}
\usepackage[utf8]{inputenc}
\usepackage[T1]{fontenc}
\usepackage{graphicx}
\usepackage{longtable}
\usepackage{wrapfig}
\usepackage{rotating}
\usepackage[normalem]{ulem}
\usepackage{amsmath}
\usepackage{amssymb}
\usepackage{capt-of}
\usepackage{hyperref}
\usepackage{fancyhdr}
\usepackage[top=25mm, left=25mm, right=25mm]{geometry}
\usepackage{longtable}
\usepackage{tabularx}
\usepackage{amsmath}
\fancyhead[R]{}
\setlength\headheight{43.0pt}
\makeatletter
\newcommand{\kc}[2]{\makebox[1.2cm][c]{\rule[-0.55cm]{0pt}{1.2cm}\in@{,#1,}{,#2,}\ifin@ 1\else 0\fi}}
\makeatother
\newcommand{\kmapthree}[1]{\begin{center}\begin{tabular}{c|*{4}{c|}}\multicolumn{1}{c}{$A\backslash BC$}&\multicolumn{1}{c}{00}&\multicolumn{1}{c}{01}&\multicolumn{1}{c}{11}&\multicolumn{1}{c}{10}\\\cline{2-5} 0&\kc{0}{#1}&\kc{1}{#1}&\kc{3}{#1}&\kc{2}{#1}\\\cline{2-5} 1&\kc{4}{#1}&\kc{5}{#1}&\kc{7}{#1}&\kc{6}{#1}\\\cline{2-5}\end{tabular}\end{center}}
\newcommand{\kmapfour}[1]{\begin{center}\begin{tabular}{c|*{4}{c|}}\multicolumn{1}{c}{$AB\backslash CD$}&\multicolumn{1}{c}{00}&\multicolumn{1}{c}{01}&\multicolumn{1}{c}{11}&\multicolumn{1}{c}{10}\\\cline{2-5} 00&\kc{0}{#1}&\kc{1}{#1}&\kc{3}{#1}&\kc{2}{#1}\\\cline{2-5} 01&\kc{4}{#1}&\kc{5}{#1}&\kc{7}{#1}&\kc{6}{#1}\\\cline{2-5} 11&\kc{12}{#1}&\kc{13}{#1}&\kc{15}{#1}&\kc{14}{#1}\\\cline{2-5} 10&\kc{8}{#1}&\kc{9}{#1}&\kc{11}{#1}&\kc{10}{#1}\\\cline{2-5}\end{tabular}\end{center}}
\usepackage[T1]{fontenc}
\usepackage[utf8]{inputenc}
\usepackage[spanish, ]{babel}
\usepackage[backend=biber]{biblatex}
\addbibresource{/mnt/d/entea/Lectures/iccd332ArqComp/bibliography/bibliography.bib}
\author{COLOQUE SU NOMBRE}
\date{2026-11-04}
\title{Tarea de Ejercicios de Circuitos Digitales Combinatorios}
\hypersetup{
 pdfauthor={COLOQUE SU NOMBRE},
 pdftitle={Tarea de Ejercicios de Circuitos Digitales Combinatorios},
 pdfkeywords={},
 pdfsubject={},
 pdfcreator={Emacs 27.1 (Org mode 9.7.5)}, 
 pdflang={Español}}
\begin{document}

\maketitle
\fancyhead[C]{\includegraphics[scale=0.05]{.././images/logoEPN.jpg}\\
ESCUELA POLITÉCNICA NACIONAL\\FACULTAD DE INGENIERÍA DE SISTEMAS\\
ARQUITECTURA DE COMPUTADORES}
\thispagestyle{fancy}

\section{Objetivos de Aprendizaje}
\label{sec:org2fa20b7}
\begin{itemize}
\item Obtener funciones booleanas en forma SOP y POS a partir de tablas de verdad.
\item Simplificar funciones booleanas con álgebra de Boole y mapas de Karnaugh.
\item Diseñar circuitos combinacionales: sumador binario y lógica de control.
\item Aplicar el complemento a 2 y el bit de signo para comparar magnitudes
(modelado de un ascensor).
\item Ejercitar el uso de Emacs Org-mode, \LaTeX{} y la exportación a PDF.
\end{itemize}

\section{Instrucciones}
\label{sec:orgce12107}
\begin{enumerate}
\item Escriba todas las expresiones matemáticas en \LaTeX{} en los espacios
provistos (reemplace los símbolos \texttt{???} o complete los renglones vacíos).
\item En los desarrollos algebraicos, cada línea del entorno \texttt{align*} es un
paso de la reducción. Debajo, liste las \textbf{propiedades del álgebra de
Boole} usadas.
\item En los ejercicios de simplificación algebraica debe \emph{además}
simplificar con \textbf{mapa de Karnaugh} y comparar ambos resultados.
\item Los circuitos se entregan como captura de simulación (Matlab/Logisim)
insertada con:
 \texttt{[[file:images/mi\_circuito.png]]}.
\item Suba al aula virtual su archivo \texttt{.org} y el \texttt{.pdf} generado.
\end{enumerate}

\section{Guía Emacs: \LaTeX{} y mapas de Karnaugh}
\label{sec:orga946249}
\subsection{Ver las fórmulas dentro de Emacs}
\label{sec:org4e9ce28}
Ubique el cursor sobre una fórmula y ejecute \texttt{C-c C-x C-l}
(\texttt{org-latex-preview}) para ver su vista previa sin generar el PDF.
Ejecute el mismo comando para ocultarla. Para generar el PDF use
\texttt{M-x org-latex-export-to-pdf}.

\subsection{Símbolos útiles de álgebra de Boole}
\label{sec:org72a22b7}
\begin{center}
\begin{tabular}{lll}
Operación & Código \LaTeX{} & Resultado\\[0pt]
\hline
NOT (negación) & \texttt{\textbackslash{}bar\{A\}} & \(\bar{A}\)\\[0pt]
NOT de expresión & \texttt{\textbackslash{}overline\{A+B\}} & \(\overline{A+B}\)\\[0pt]
XOR & \texttt{A \textbackslash{}oplus B} & \(A \oplus B\)\\[0pt]
AND / OR & \texttt{AB}, \texttt{A+B} & \(AB\), \(A+B\)\\[0pt]
\end{tabular}
\end{center}

\subsection{Desarrollos alineados con \texttt{align*}}
\label{sec:org2a26ac7}
Para alinear los pasos de una reducción respecto del signo \texttt{=} use \texttt{\&=}
al inicio de cada línea y \texttt{\textbackslash{}\textbackslash{}} al final (excepto la última). Anote a la
derecha la propiedad con \texttt{\textbackslash{}text\{\}} si lo desea:

\begin{verbatim}
\begin{align*}
  F &= ABD + \bar{A}B\bar{C}D \\
    &= BD\,(A + \bar{A}\bar{C}) \\
    &= \dots
\end{align*}
\end{verbatim}

Resultado:

\begin{align*}
  F &= ABD + \bar{A}B\bar{C}D \\
    &= BD\,(A + \bar{A}\bar{C}) \\
    &= \dots
\end{align*}

\subsection{Dibujar un mapa K con \LaTeX{}}
\label{sec:org7901487}
Esta tarea ya define dos comandos en el encabezado que dibujan mapas
con las columnas en \textbf{código Gray} (00, 01, 11, 10):

\begin{itemize}
\item \texttt{\textbackslash{}kmapthree\{lista\}}: mapa de 3 variables (\(A\) filas; \(BC\) columnas).
\item \texttt{\textbackslash{}kmapfour\{lista\}}: mapa de 4 variables (\(AB\) filas; \(CD\) columnas).
\end{itemize}

Se insertan dentro de un bloque de exportación (ya están colocados en
cada ejercicio).

\textbf{Cómo llenar el mapa.} En cada ejercicio encontrará el comando con las
llaves vacías, por ejemplo \texttt{\textbackslash{}kmapthree\{\}}. Usted debe:

\begin{enumerate}
\item Determinar los \textbf{min-términos} donde \(F=1\) (en una tabla de verdad,
el número de fila en decimal; para una expresión, evalúela o
expándala hasta min-términos). Con \(A\) como bit más significativo,
\(m_6=110_2\) corresponde a \(A=1,B=1,C=0\).
\item Escribir esos números \textbf{dentro de las llaves}, separados por comas y
\textbf{sin espacios}.
\item Exportar o previsualizar: el comando coloca un 1 en las celdas de la
lista y un 0 en las demás, y ordena las columnas en \textbf{código Gray}
(00, 01, 11, 10).
\end{enumerate}

Ejemplo (función distinta a las de la tarea): para
\(F=\sum m(1,3,4)\) se escribe

\begin{verbatim}
#+begin_export latex
\kmapthree{1,3,4}
#+end_export
\end{verbatim}

y se obtiene

\kmapthree{1,3,4}

Con 4 variables funciona igual; por ejemplo \texttt{\textbackslash{}kmapfour\{0,2,8,10\}}
dibuja \(F=\sum m(0,2,8,10)\).

Las agrupaciones no se dibujan automáticamente (no se generalizan de
forma eficiente): bajo el mapa liste los grupos formados y el término
que genera cada uno, y finalmente la función simplificada con
\texttt{align*}.

\emph{Opcional:} el paquete \texttt{karnaugh-map} (requiere \texttt{tikz}) dibuja mapas
con agrupaciones coloreadas. Añádalo con
\texttt{\#+latex\_header: \textbackslash{}usepackage\{karnaugh-map\}} y use el entorno
\texttt{karnaugh-map} (consulte la documentación del paquete con \texttt{texdoc karnaugh-map}).

\section{Parte A. Álgebra de Boole}
\label{sec:orgb653e82}
\subsection{Ejercicio 1: Compuerta XOR en SOP}
\label{sec:orge8e8ba0}
A partir de la tabla de verdad de la compuerta OR exclusiva de dos
entradas obtenga la función booleana como suma de productos (SOP,
min-términos). Complete la tabla y la expresión.

\begin{center}
\begin{tabular}{rrl}
A & B & \(A \oplus B\)\\[0pt]
\hline
0 & 0 & \\[0pt]
0 & 1 & \\[0pt]
1 & 0 & \\[0pt]
1 & 1 & \\[0pt]
\end{tabular}
\end{center}

\textbf{Respuesta (SOP):}

$$A \oplus B = ???$$

\subsection{Ejercicio 2: Compuerta XOR en POS}
\label{sec:org789d16e}
Para el ejercicio anterior obtenga la representación como producto de
sumas (POS, max-términos).

\textbf{Respuesta (POS):}

$$A \oplus B = ???$$

\subsection{Ejercicio 3: Implementación con NAND}
\label{sec:org1a382ca}
¿Puede representar la XOR usando solamente compuertas NAND? Aplique las
leyes de De Morgan, recordando que una NAND es una AND negada
(\(\overline{AB}\)). Desarrolle los pasos hasta obtener una expresión
únicamente con NAND y dibuje el circuito.

\textbf{Respuesta:}

\begin{align*}
  A \oplus B &= \bar{A}B + A\bar{B} \\
             &= \\
             &=
\end{align*}

\textbf{Propiedades del álgebra de Boole usadas:}
\begin{enumerate}
\item Propiedad \ldots{}
\item Propiedad \ldots{}
\item Propiedad \ldots{}
\end{enumerate}

\textbf{Circuito (captura o descripción):}

\subsection{Ejercicio 4: Simplificación algebraica I y mapa K}
\label{sec:org266823e}
Simplifique \(F = ACD+\bar{A}BCD\).

\textbf{a) Con álgebra de Boole:}

\begin{align*}
  F &= ACD+\bar{A}BCD \\
    &= \\
    &=
\end{align*}

\textbf{Propiedades del álgebra de Boole usadas} (indique el paso en que se usa cada una):
\begin{enumerate}
\item Propiedad \ldots{}
\item Propiedad \ldots{}
\item Propiedad \ldots{}
\end{enumerate}

\textbf{b) Con mapa de Karnaugh (4 variables).} Escriba los min-términos
donde \(F=1\) dentro de las llaves de \texttt{\textbackslash{}kmapfour\{\}}:

\kmapfour{}

Min-términos: \(F=\sum m(\ldots)\)

Grupos formados y término que genera cada uno:

\begin{align*}
  G_1 &= \\
  G_2 &=
\end{align*}

\begin{align*}
  F_{\text{Karnaugh}} &=
\end{align*}

\textbf{c) Compare:} ¿coinciden los resultados de (a) y (b)? Comente.

\textbf{Respuesta:}

\vspace{1.5cm}

\subsection{Ejercicio 5: Simplificación algebraica II y mapa K}
\label{sec:org40cad12}
Simplifique \(F=ABC+A\bar{B}\,\overline{\bar{A}\bar{C}}\).

\textbf{a) Con álgebra de Boole:}

\begin{align*}
  F &= ABC+A\bar{B}\,\overline{\bar{A}\bar{C}} \\
    &= \\
    &=
\end{align*}

\textbf{Propiedades del álgebra de Boole usadas:}
\begin{enumerate}
\item Propiedad \ldots{}
\item Propiedad \ldots{}
\item Propiedad \ldots{}
\end{enumerate}

\textbf{b) Con mapa de Karnaugh (3 variables).} Obtenga los min-términos de
la expresión original y escríbalos dentro de las llaves de \texttt{\textbackslash{}kmapthree\{\}}:

\kmapthree{}

Min-términos: \(F=\sum m(\ldots)\)

Grupos formados y término que genera cada uno:

\begin{align*}
  G_1 &= \\
  G_2 &=
\end{align*}

\begin{align*}
  F_{\text{Karnaugh}} &=
\end{align*}

\textbf{c) Compare:} ¿coinciden los resultados de (a) y (b)? Comente.

\textbf{Respuesta:}

\vspace{1.5cm}

\section{Parte B. Tablas de verdad y mapas de Karnaugh}
\label{sec:orgcbd1e57}
\subsection{Ejercicio 6: SOP desde tabla de verdad}
\label{sec:org57fa48d}
A partir de la tabla de verdad de tres variables obtenga la
representación en SOP.

\begin{table}[htbp]
\caption{\label{tab:org1e03ec0}Ejercicio de tres variables}
\centering
\begin{tabular}{rrrr}
A & B & C & F\\[0pt]
\hline
0 & 0 & 0 & 0\\[0pt]
0 & 0 & 1 & 0\\[0pt]
0 & 1 & 0 & 1\\[0pt]
0 & 1 & 1 & 1\\[0pt]
1 & 0 & 0 & 0\\[0pt]
1 & 0 & 1 & 0\\[0pt]
1 & 1 & 0 & 1\\[0pt]
1 & 1 & 1 & 0\\[0pt]
\end{tabular}
\end{table}

\textbf{Respuesta:}

\begin{align*}
  F(A,B,C) &= \sum m(\ldots) \\
           &=
\end{align*}

\subsection{Ejercicio 7: Simplificación con mapa de Karnaugh}
\label{sec:org49ffa02}
Usando mapas de Karnaugh simplifique la función de la tabla del
Ejercicio 6 y dibuje el circuito resultante.

Escriba los min-términos de \(F\) (Ej. 6) dentro de las llaves:

\kmapthree{}

Grupos formados:

\begin{align*}
  G_1 &= \\
  G_2 &=
\end{align*}

\begin{align*}
  F_{\text{simplificada}} &=
\end{align*}

\textbf{Circuito:}

\subsection{Ejercicio 8: Karnaugh y simulación}
\label{sec:org0e8cae6}
Dada la siguiente tabla de verdad, simplifique utilizando un mapa de
Karnaugh, obtenga el circuito combinacional y simule el resultado.

\begin{table}[htbp]
\caption{\label{tab:org320d085}Ejercicio 8}
\centering
\begin{tabular}{rrrr}
A & B & C & F\\[0pt]
\hline
0 & 0 & 0 & 1\\[0pt]
0 & 0 & 1 & 1\\[0pt]
0 & 1 & 0 & 0\\[0pt]
0 & 1 & 1 & 1\\[0pt]
1 & 0 & 0 & 0\\[0pt]
1 & 0 & 1 & 1\\[0pt]
1 & 1 & 0 & 0\\[0pt]
1 & 1 & 1 & 1\\[0pt]
\end{tabular}
\end{table}

Escriba los min-términos de \(F\) dentro de las llaves:

\kmapthree{}

Grupos formados:

\begin{align*}
  G_1 &= \\
  G_2 &=
\end{align*}

\begin{align*}
  F_{\text{simplificada}} &=
\end{align*}

\textbf{Captura de la simulación:}

\begin{verbatim}
[[file:images/ejercicio8.png]]
\end{verbatim}

\section{Parte C. Sumador binario}
\label{sec:orgdfe4b0f}
\subsection{Ejercicio 9: Sumador completo en cascada}
\label{sec:org336fd1f}
Obtenga un circuito combinacional que realice la suma binaria. Considere
que el sumador se resuelve integrando en cascada varias etapas que
suman los dígitos de los sumandos y el acarreo, que inicia en 0.
Consulte \autocite{stallings2006esp}, páginas PDF 777 a 780 y
\autocite{stallings2022eng}, páginas 435 a 438.

\begin{enumerate}
\item Complete la tabla de verdad de una etapa (sumador completo).

\begin{center}
\begin{tabular}{rrrll}
\(A_i\) & \(B_i\) & \(C_i\) & \(S_i\) & \(C_{i+1}\)\\[0pt]
\hline
0 & 0 & 0 &  & \\[0pt]
0 & 0 & 1 &  & \\[0pt]
0 & 1 & 0 &  & \\[0pt]
0 & 1 & 1 &  & \\[0pt]
1 & 0 & 0 &  & \\[0pt]
1 & 0 & 1 &  & \\[0pt]
1 & 1 & 0 &  & \\[0pt]
1 & 1 & 1 &  & \\[0pt]
\end{tabular}
\end{center}

\item Deduzca las expresiones booleanas reducidas (use mapas K de 3 variables).

\begin{align*}
  S_i     &= \\
  C_{i+1} &=
\end{align*}

\item Implemente el circuito en Matlab (o Logisim), simule e inserte la captura.

\begin{verbatim}
   [[file:images/sumador.png]]
\end{verbatim}

\item Extienda el diseño a un sumador de 4 bits en cascada. Dibuje el
diagrama de bloques.
\end{enumerate}

\section{Parte D. Modelado de un ascensor}
\label{sec:org5f797b2}
\subsection{Ejercicio 10: Lógica de control de un ascensor de 10 pisos}
\label{sec:org2572d94}
Se desea modelar la lógica combinacional de un ascensor de \textbf{10 pisos}
(pisos 0 a 9). El sistema recibe:

\begin{itemize}
\item \(P_{act}\): piso actual (4 bits, \(P_3P_2P_1P_0\), valores 0 a 9).
\item \(P_{dest}\): piso destino (4 bits, \(D_3D_2D_1D_0\), valores 0 a 9).
\end{itemize}

y debe generar tres salidas: \textbf{SUBIR}, \textbf{BAJAR} y \textbf{DETENER}.

\textbf{Idea de diseño.} La diferencia \(\Delta = P_{dest}-P_{act}\) se calcula
con un sumador y el complemento a 2:
\[
\Delta = P_{dest} + \overline{P_{act}} + 1 = P_{dest} + C_2(P_{act}).
\]
Como \(\Delta \in [-9,9]\), se requieren \textbf{5 bits} (extensión de signo de
los operandos a 5 bits). El \textbf{bit de signo} \(\Delta_4\) indica el sentido:

\begin{itemize}
\item \(\Delta_4 = 1\) (negativo): el destino está por debajo, el ascensor \textbf{baja}.
\item \(\Delta_4 = 0\) y \(\Delta \neq 0\): el destino está por encima, \textbf{sube}.
\item \(\Delta = 0\): el ascensor ya está en el piso, \textbf{se detiene}.
\end{itemize}

\subsubsection{10.1 Representación}
\label{sec:org456e05d}
Complete la tabla con la representación en 5 bits (complemento a 2) de
\(-P_{act}\) y de \(\Delta\).

\begin{center}
\begin{tabular}{rrllll}
\(P_{act}\) & \(P_{dest}\) & \(C_2(P_{act})\) (5 bits) & \(\Delta\) (5 bits) & \(\Delta_4\) & Acción\\[0pt]
\hline
3 & 7 &  &  &  & \\[0pt]
8 & 2 &  &  &  & \\[0pt]
5 & 5 &  &  &  & \\[0pt]
0 & 9 &  &  &  & \\[0pt]
9 & 0 &  &  &  & \\[0pt]
\end{tabular}
\end{center}

\subsubsection{10.2 Complemento a 2 y sumador}
\label{sec:org248029a}
Dibuje el diagrama de bloques del circuito que calcula \(\Delta\):
inversores sobre \(P_{act}\), sumador de 5 bits en cascada (Ejercicio 9) y
acarreo de entrada \(C_0=1\).

\subsubsection{10.3 Detector de cero y salidas}
\label{sec:org9723b55}
Obtenga las expresiones de las salidas en función de los bits de
\(\Delta\) (\(\Delta_4\ldots\Delta_0\)).

\begin{align*}
  \text{CERO}    &= \\
  \text{BAJAR}   &= \\
  \text{SUBIR}   &= \\
  \text{DETENER} &=
\end{align*}

\subsubsection{10.4 Circuito completo}
\label{sec:org0226364}
Dibuje el circuito combinacional completo (complemento a 2, sumador,
detector de cero y salidas) y simule al menos los cinco casos de la
tabla de 10.1. Inserte la captura.

\subsubsection{10.5 Análisis}
\label{sec:org45c2269}
a) ¿Por qué es insuficiente usar solo 4 bits para \(\Delta\)? Dé un
   ejemplo con pisos concretos.

\textbf{Respuesta:}

\vspace{1.5cm}

b) Las combinaciones con valores 10 a 15 no son válidas. Proponga cómo
   tratarlas (condiciones \emph{don't care} o señal de error) y su efecto en
   la simplificación.

\textbf{Respuesta:}

\vspace{1.5cm}

c) ¿Qué elementos secuenciales harían falta para que el ascensor
   \emph{memorice} las solicitudes y el piso actual?

\textbf{Respuesta:}

\vspace{1.5cm}

\section{Rúbrica de evaluación}
\label{sec:orgde0b760}
\begin{center}
\begin{tabular}{lr}
Criterio & Puntaje máximo\\[0pt]
\hline
Parte A: Álgebra de Boole y mapas K (Ej. 1-5) & 30\\[0pt]
Parte B: Tablas de verdad y Karnaugh (Ej. 6-8) & 20\\[0pt]
Parte C: Sumador binario (Ej. 9) & 15\\[0pt]
Parte D: Ascensor (Ej. 10) & 30\\[0pt]
Presentación, uso de Org-mode/\LaTeX{} y PDF & 5\\[0pt]
TOTAL & 100\\[0pt]
\end{tabular}
\end{center}

\section{Verificación de Entregables [0\%]:}
\label{sec:org97c0854}
\begin{itemize}
\item[{$\square$}] Todos los ejercicios resueltos y justificados en \LaTeX{}
\item[{$\square$}] Propiedades del álgebra de Boole listadas en cada simplificación
\item[{$\square$}] Capturas de simulaciones insertadas
\item[{$\square$}] Revisión de ortografía con \texttt{ispell-buffer}
\item[{$\square$}] Generación de PDF \texttt{M-x org-latex-export-to-pdf}
\item[{$\square$}] Subir al aula virtual \texttt{.org} y \texttt{.pdf}
\end{itemize}

\subsection{Problemas con la Tarea/logísticos:}
\label{sec:org2db2d0f}
\begin{itemize}
\item Ejercicio \(X_1\) no pudo completarse debido a
\end{itemize}

\section{Referencias}
\label{sec:orgbd05b6e}
\printbibliography
\end{document}
